\section{Related Work}

Target-supervised calibration uses labeled target examples to adjust calibration parameters, with prior work differing in when these examples become available. Online calibration updates parameters during operation as labeled target examples arrive sequentially~\cite{chen2024online}. This approach achieves a mean calibration error of 0.08 against a baseline of 0.12 across three corruptions. Pre-evaluation calibration instead fits a scalar temperature using 500 labeled target examples while keeping classifier parameters fixed~\cite{rao2025frozen}. Under this setting, the mean calibration error is 0.09 against a baseline of 0.13 on two fog datasets.

Our study follows the same pre-evaluation setting~\cite{rao2025frozen}, fitting a scalar temperature with 500 labeled target images and fixed classifier parameters. This places target supervision before evaluation, in contrast to online calibration with sequential updates~\cite{chen2024online}. Building on this established calibration procedure, we examine the measured tradeoff between calibration quality and latency across three fixed target-domain test sets. The contribution therefore concerns how this tradeoff varies across domains within a shared target-supervised calibration setting.
